% !TEX root = ../main.tex
\chapter{Repos, Postponing Settlement}\label{ch:repo}
\begin{paracol}{2}

\begin{sources}
\begin{tabularx}{\linewidth}{@{}l l L@{}}
\toprule
\textbf{Segment} & \textbf{Length} & \textbf{Content}\\
\midrule
\seg{7}{1}{L7.1 FT: The Impact of QE3} & 2:45 & Second thoughts on QE3: buying time, and hedge funds betting against it.\\
\seg{7}{2}{L7.2 Money Market Interest Rate Patterns} & 3:34 & LIBOR, Fed funds, repo: the usual order, and its reversal.\\
\seg{7}{3}{L7.3 What is Repo?} & 3:19 & A collateralized loan organised as a sale and repurchase.\\
\seg{7}{4}{L7.4 Repo in Balance Sheets} & 7:34 & Repo and reverse repo around a security dealer; the confusing language.\\
\seg{7}{5}{L7.5 Comparison with Fed Funds} & 5:09 & Repo as an open, dealer-made market.\\
\seg{7}{6}{L7.6 Legal Construction of Repo} & 9:46 & Symmetry, margin, haircut, repo interest (Stigum's example).\\
\seg{7}{7}{L7.7 Security Dealers Balance Sheet} & 11:07 & Primary dealer statistics; dealers as banks; haircuts and liquidity spirals.\\
\seg{7}{8}{L7.8 Repo, Modern Finance, and the Fed} & 8:49 & Repo as money; Fed open market operations as repos with dealers.\\
\seg{7}{9}{L7.9 Interest Rate Spreads: Before the Crisis} & 5:33 & Why repo $<$ Fed funds: Bagehot's market rate and bank rate.\\
\seg{7}{10}{L7.10 Interest Rate Spreads: After the Crisis} & 8:07 & Repo $>$ Fed funds: stress in dollar funding; reading spreads as clues.\\
\bottomrule
\end{tabularx}

\smallskip
\textbf{Files.} Transcripts \file{materials/transcripts/L07-P*.txt}. Segment~5 has no YouTube captions: its transcript
\file{L07-P05.txt} is extracted from the original \file{transcript-lecture-7to12.txt}, so its times may be off by a few seconds.
\textbf{Reading.} Stigum, \emph{The Money Market}, chapter on repo (pages 533--534 are quoted in class).
\end{sources}

\section*{Motivation}
The Fed funds market (lecture~6) lets banks meet at the end of the day and push settlement to tomorrow. The repo market does the same, but for
almost everybody and against collateral: it is the core funding market of the security dealers, and through them of modern finance. The lecture
explains what repo is (\cref{sec:rp-what}), how it appears in balance sheets (\cref{sec:rp-bs}), how it is priced (\cref{sec:rp-legal}), why
dealers that fund themselves this way are banks (\cref{sec:rp-dealers}), how the Fed uses it (\cref{sec:rp-fed}), and what the spread between repo
and Fed funds tells us about the system (\cref{sec:rp-spreads}).

% ------------------------------------------------------------------
\section{FT: the impact of QE3}\label{sec:rp-ft}

\begin{casestudy}[FT, September 2012: second thoughts on QE3]
John Plender, ``Central bank firepower risks creating false sense of security'': QE3 puts a floor under certain security prices but does nothing
fundamental; if the problem is fiscal, central banks only buy time, and if nobody uses the time, nothing has been done. His last paragraph:
``throughout history monetary experiments have shown a nasty tendency to blow up\dots\ this experiment is vast; it must not breed complacency''
\ts{7}{1}{0:06}. On the same page, ``Hedge fund sceptics take on the bulls over QE infinity'': some funds argue that QE1
worked well, QE2 less, QE3 may not even lift asset prices, and trade against those who expect a boom \ts{7}{1}{1:08}.
\end{casestudy}
If private investors trade against it, QE buys even less time \ts{7}{1}{2:09}. Central banks act not because they are optimistic, but because
nobody else is doing anything; and in repo rates we can see symptoms of the stresses in the system \ts{7}{1}{2:09}.

% ------------------------------------------------------------------
\section{Money market interest rate patterns}\label{sec:rp-patterns}

In the FT: US dollar LIBOR overnight 0.1509\%, US overnight repo 0.31\%, Fed funds effective 0.16\% \ts{7}{2}{0:07}. Before the crisis
the typical order was
\[
  \text{repo} \;<\; \text{Fed funds} \;<\; \text{LIBOR (Eurodollar)} ,
\]
and Mehrling had a whole story for it, from a paper he wrote before the crisis \ts{7}{2}{0:56}. Now repo is \emph{above} Fed funds:
odd, since repo is secured and Fed funds are not \ts{7}{2}{1:38}.

\begin{nfigure}
\begin{tikzpicture}[font=\scriptsize,x=1cm,y=1cm]
  \draw[-{Stealth[length=3pt]}] (0,0) -- (8,0) node[right]{time (mid-2005 to mid-2006)};
  \draw[-{Stealth[length=3pt]}] (0,-1.8) -- (0,1.6) node[above]{spread vs target};
  \node[left] at (0,0) {0};
  \draw[thick,intcol] plot[smooth] coordinates {(0,0.45) (1,0.5) (1.5,-0.6) (1.7,0.45) (3,0.5) (3.5,-0.5) (3.7,0.4) (5,0.45) (5.5,-0.7) (5.7,0.4) (7.5,0.5)};
  \draw[thick,thmcol] plot[smooth] coordinates {(0,-0.4) (1,-0.35) (1.5,-1.5) (1.7,-0.4) (3,-0.35) (3.5,-1.3) (3.7,-0.4) (5,-0.35) (5.5,-1.6) (5.7,-0.4) (7.5,-0.35)};
  \node[intcol,right] at (7.5,0.5) {target $-$ repo $>0$};
  \node[thmcol,right] at (7.5,-0.35) {target $-$ Eurodollar $<0$};
\end{tikzpicture}
\caption{Schematic of the pre-crisis chart shown in class \ts{7}{2}{1:58}: Fed funds \emph{target} minus repo was typically positive, target
minus Eurodollar typically negative. The downward spikes occur before each step of the 2005--2006 tightening, as market rates rise in anticipation
\ts{7}{2}{2:39}. (Shape illustrative.)}\label{fig:rp-spreads}
\end{nfigure}

% ------------------------------------------------------------------
\section{What is repo?}\label{sec:rp-what}

\begin{definition}[New concepts: repurchase agreement]\label{def:rp}
A \emph{repurchase agreement} (``repo'', ``RP'') is a loan, overnight or for a term, collateralized by a security, organised as a
\emph{sale and repurchase}: the borrower of money sells a security today and promises to buy it back tomorrow at a fixed price; the securities go
from the borrower to the lender and back, the money from the lender to the borrower and back \ts{7}{3}{0:08}.
(Stigum draws the ``first leg'' and the ``second leg'' of the repo \ts{7}{3}{1:34}.)
\end{definition}

The securities are collateral in a strong sense: if the borrower does not repay, the lender \emph{already holds} the securities and can sell them;
there is no need to depend on the borrower's credit \ts{7}{3}{1:55}. And symmetrically: if the lender does not return the securities, the
borrower need not repay the money. Such a \emph{fail} turns into a free overnight continuation of the loan until the securities are delivered
\ts{7}{3}{2:37}. The money collateralizes the loan of securities, and the securities collateralize the loan of money.

\begin{nfigure}
\begin{tikzpicture}[font=\small,
  ag/.style={draw,rounded corners=2pt,minimum width=2.6cm,minimum height=0.8cm,align=center}]
  \node[ag,fill=defcol!10] (b) at (0,0) {borrower of money\\(cash taker)};
  \node[ag,fill=intcol!10] (l) at (6,0) {lender of money\\(cash provider)};
  \draw[-{Stealth[length=4pt]},thick,keycol] ([yshift=4pt]b.east) -- node[above,font=\scriptsize]{today: securities} ([yshift=4pt]l.west);
  \draw[-{Stealth[length=4pt]},thick,fincol] ([yshift=-4pt]l.west) -- node[below,font=\scriptsize]{today: money} ([yshift=-4pt]b.east);
  \draw[-{Stealth[length=4pt]},dashed,keycol] (l.south) to[out=-110,in=-70] node[below,font=\scriptsize]{tomorrow: securities back} (b.south);
  \draw[-{Stealth[length=4pt]},dashed,fincol] (b.north) to[out=70,in=110] node[above,font=\scriptsize]{tomorrow: money $+$ interest} (l.north);
\end{tikzpicture}
\caption{The two legs of a repo \ts{7}{3}{1:14}: a sale of securities today and their repurchase tomorrow at a slightly higher price.}\label{fig:rp-legs}
\end{nfigure}

% ------------------------------------------------------------------
\section{Repo in balance sheets}\label{sec:rp-bs}

There are trillions of these overnight loans \ts{7}{4}{0:07}. The language is not standardised: ``doing repo'', ``reversing in securities'', ``repoing
securities''; and, Mehrling warns, the same phrase can mean opposite sides of the trade in New York and in London. Balance sheets keep you straight \ts{7}{4}{0:27}.
Two rules:
\begin{enumerate}
  \item almost every repo has a \emph{security dealer} on one side \ts{7}{4}{0:48};
  \item ``follow the money'': a repo booked as a liability means money borrowed, as an asset money lent; money is better than securities, so its side tells
        you where to book \ts{7}{4}{1:54}.
\end{enumerate}

\begin{taccts}
\tacct[2.0cm]{Pension fund}{Repo loan to dealer}{}\quad
\tacct[2.0cm]{Security dealer}{Reverse repo (loan to bank)}{Repo (from pension fund)}\quad
\tacct[2.0cm]{Bank}{Securities}{Repo (from dealer)}
\end{taccts}

\begin{itemize}
  \item \textbf{Pension fund}: lends money overnight, against a Treasury, for interest. For it this is practically a deposit, a way to hold cash, and
        better than a deposit, because large deposits are not insured: if the bank fails the cash is gone, while in repo it holds a Treasury
        \ts{7}{4}{2:15}.
  \item \textbf{Dealer lending money}: takes in securities; Stigum calls it a \emph{reverse} \ts{7}{4}{3:16}. Securities move one way across the
        diagram, money the other \ts{7}{4}{3:57}.
  \item \textbf{Bank}: holds securities and funds them by repoing them out overnight, instead of with deposits: repo as the liability that funds an asset
        \ts{7}{4}{4:17}.
\end{itemize}
The dealer's ``repo'' and ``reverse'' are the same instrument seen from the two sides; they get different names because the dealer is in the middle,
where one is a liability and the other an asset \ts{7}{4}{5:40}. To add to the confusion, the Fed uses the words the other way round: when it
borrows from dealers it calls it a ``reverse repo'', when it lends (open market operations) a ``repo'' \ts{7}{4}{6:21}. Always translate
into balance sheets: who is borrowing money? \ts{7}{4}{7:23}

\begin{history}[Repo before the Fed]
Repo is ``very ancient'': in the United States it existed before the Fed, as the way banks traded with each other before the Fed funds market. With a
correspondent you could rely on the relationship; with a stranger you wanted collateral: lend against a bond, and you only need to check the bond, not the
borrower. Mehrling links the rise of bond ratings (Moody's) to this need to know how much to lend against which bond \ts{7}{4}{4:38}.
(John Moody published his first bond ratings in 1909.)\par
\emph{References:} R.~Sylla, ``An historical primer on the business of credit rating'', in R.~Levich et al.\ (eds.), \emph{Ratings, Rating Agencies and the
Global Financial System}, Kluwer, 2002.
\end{history}

\paragraph{Comparison with Fed funds.} The repo market is like the Fed funds market: it lets deficit and surplus agents meet and push their liabilities and
assets to another day, with a dealer in the middle \ts{7}{5}{0:03}. Two differences \ts{7}{5}{1:25}:
\begin{enumerate}
  \item \textbf{Openness}: anyone with a Treasury bond can do repo, without being a bank or a member of the Federal Reserve System: it is ``much more
        democratic'' \ts{7}{5}{1:33};
  \item \textbf{Dealing versus brokering}: in Fed funds there is little dealing, it is mostly brokered; in repo it is almost all dealing, on the balance
        sheets of security dealers \ts{7}{5}{0:52}.
\end{enumerate}
Why a dealer market? Repo is heterogeneous: any security can be repoed (corporate bonds, mortgage-backed securities, \dots) and anyone can participate.
The \emph{name of the dealer}, a primary dealer, a big bank, on the other side of every trade, is what makes the market homogeneous: you only worry about
the dealer \ts{7}{5}{2:12}. There is even \emph{general collateral} repo, in which the dealer may substitute one
security for another of the same class, to keep things smooth; the dealer is a liability on both sides, to return securities and to return money
\ts{7}{5}{3:53}. Repo defaults usually happen not because someone is insolvent but because the securities were sold on and not returned
\ts{7}{5}{4:42}.

% ------------------------------------------------------------------
\section{The legal construction of repo: haircut and interest}\label{sec:rp-legal}

Repo is like a collateralized loan, but much more symmetric \ts{7}{6}{0:07}. With a mortgage, the bank has a claim on the house, but foreclosure takes months
and legal fees, and nobody intends to use it \ts{7}{6}{0:28}. A repo is an overnight loan of \$10 million: the lender \emph{owns} the security, can sell it or
re-pledge it (\emph{rehypothecate}) and trade with it all day, and is only liable for returning it the next day; the borrower can spend the money as it likes
\ts{7}{6}{0:48}.

\paragraph{Who gives margin to whom?} The lender of money fears that the security falls below the value of the loan, which gives the borrower an incentive not
to repay; the lender of securities fears the opposite \ts{7}{6}{1:51}. Stigum's answer: it is always the \emph{lender of money} who gets margin
\ts{7}{6}{2:31}. ``This is part of the hierarchy of money and credit'': money is the better asset, and its holder can demand extra security; and the demand for
liquidity is the short side of the market, people pay up to get money, not to get collateral \ts{7}{6}{2:31}. (In swaps, by contrast,
typically both sides post margin \ts{7}{6}{9:09}.)

\begin{definition}[New concepts: haircut, repo rate]
The \emph{haircut} is the percentage by which the amount lent is smaller than the market value of the collateral; it protects the lender of money against a fall
in the collateral's price. The \emph{repo rate} is the interest rate on the loan, quoted on a simple-interest, actual/360 basis.
\end{definition}

\paragraph{Stigum's example} \ts{7}{6}{2:52}--\ts{7}{6}{7:47}. A 10-year Treasury bond trades a bit above par (its coupon is higher than the current 10-year
rate). It is offered as collateral for an overnight loan with a haircut of 2\%, so the amount lent per 100 of face value is 99.228. The repo rate is 4.92\%
(annualised). The next day the borrower repays principal plus simple interest for one day on a 360-day year:
\begin{keyformulas}
\[
  \text{repayment} = L\left(1 + r_{\text{repo}}\,\frac{d}{360}\right),\qquad L = (1-h)\,P ,
\]
with $P$ the market value of the collateral, $h$ the haircut, $d$ the number of days (1 overnight, 3 over a weekend, 7 for a week).
\end{keyformulas}
The conventions (simple, not compound interest; a 360-day year) date ``from time immemorial'': the market prices them in, and overnight they change the result
only at the fourth decimal, but traders must get them right \ts{7}{6}{6:04}. On a \$1 million bond the overnight interest
is about \$137: the cost of borrowing a million dollars overnight \ts{7}{6}{7:26}.

\begin{coursenote}[Checking the numbers]
The price of the bond is garbled in the captions (``129 30 seconds''); the other numbers are consistent with a price of about $101\tfrac{8}{32}=101.25$:
$0.98\times101.253=99.228$. Interest on \$1 million for one day at 4.92\%/360 is \$136.67, the ``137'' of the lecture; on the amount actually lent,
\$992{,}280, it is \$135.61. Bond prices were quoted in 32nds of a point (one point $=$ 1\% of face value).
\end{coursenote}

\paragraph{Haircuts make collateral homogeneous.} Different securities get different haircuts (2\% for a Treasury, 5, 10 or 50\% for a mortgage-backed
security, which in a crisis may not be accepted at all) and then the same repo rate \ts{7}{6}{4:38}. The haircut is negotiated first,
according to the volatility of the collateral's price, then the rate \ts{7}{6}{4:59}.

% ------------------------------------------------------------------
\section{The security dealer's balance sheet}\label{sec:rp-dealers}

The New York Fed publishes weekly statistics on the primary dealers \ts{7}{7}{0:07}:
\begin{itemize}
  \item \textbf{Transactions} by type of security: Treasury coupon securities, agencies (Fannie Mae, Freddie Mac), mortgage-backed, corporate
        \ts{7}{7}{0:48}.
  \item \textbf{Net positions} (inventory, hence risk exposure): about \$26 billion of Treasury bills, \$68 billion of coupon securities due in three years or
        less, negligible longer bonds \ts{7}{7}{1:08}. This is unusual: normally dealers are \emph{short} bills and \emph{long} bonds, picking up the term
        spread. With Operation Twist the Fed buys bonds from the dealers and sells them bills, sucking out their bond inventory and wiping out their short in
        bills; the dealers absorb the mismatch until they can distribute it \ts{7}{7}{1:49}. Analysts from
        Barclays, in a Reuters article, explained the high repo rate this way; Mehrling wonders why it should affect the repo rate \ts{7}{7}{2:30}.
        (Net positions say nothing about hedges with derivatives \ts{7}{7}{4:14}.)
  \item \textbf{Financing} (gross): securities in, about \$1.9 trillion, financed overnight or at term, and securities out (short positions)
        \ts{7}{7}{4:35}.
\end{itemize}
The memorandum items give the repo book, in billions \ts{7}{7}{5:38}:
\begin{taccts}
\tacct[3.0cm]{Primary dealers (\$ billion)}{Reverse repo, overnight\amt{878}\\Reverse repo, term\amt{1{,}316}}{Repo, overnight\amt{1{,}842}\\Repo, term\amt{931}}
\end{taccts}
Dealers fund their long positions in the repo market and their short positions in the reverse market \ts{7}{7}{6:44}. And look at the maturities: the big number
on the asset side is \emph{term}, on the liability side \emph{overnight}. Dealers borrow short and lend long, ``just like a bank'' (term may be a week, but it is
more than a day) \ts{7}{7}{7:25}.

\begin{proposition}[Security dealers are banks]
Overnight repo to dealers is, for institutional investors, a deposit, even better than one because secured; at \$1.8 trillion it is larger than M1. Dealers,
trading securities for profit, create an enormous volume of monetary assets: the wholesale money supply \ts{7}{7}{7:04}.
\end{proposition}

\paragraph{Haircuts and liquidity spirals.} Repo must be rolled over, and everything is ``up for grabs'' at each renewal: haircut and rate
\ts{7}{7}{9:28}. If the haircut on a security you finance with repo jumps overnight from 2\% to 10\%, you have a \emph{funding gap}: yesterday you
could borrow 98 cents on the dollar, today 90; if you cannot find the other eight, you must sell the security at a fire-sale price, which pushes prices down
further. In the crisis, haircuts on ``triple-A'' mortgage-backed securities went from 2\% to 10\% to 50\%, and the scramble for funding distorted every other price
in the world \ts{7}{7}{9:48}.

\begin{intuition}[Leverage and the haircut]
With haircut $h$, a dealer with capital $E$ can hold securities worth at most $E/h$: leverage $1/h$ (50 at $h=2\%$, 10 at $10\%$, 2 at $50\%$). A rise of the
haircut from 2\% to 10\% forces assets to fall by a factor of five for the same capital: a positive-feedback loop if the sales lower prices and raise haircuts
further.
\end{intuition}

% ------------------------------------------------------------------
\section{Repo, modern finance and the Fed}\label{sec:rp-fed}

Most money and banking textbooks do not mention repo: they are about retail money and banks, and a security dealer ``is not part of our subject''. This is where
modern finance changed the world: not being called a bank does not mean not operating like one, and producing liquid liabilities is producing money supply for
someone \ts{7}{8}{0:07}. The lender in repo may be a money market mutual fund, which holds repo and issues shares that function as deposits, to
corporations or even foreign central banks: money, directly or passed through a fund \ts{7}{8}{0:49}.

\paragraph{Open market operations are repos with dealers.} The textbook says the Fed buys or sells Treasury securities to change the money supply; in reality it
trades with the primary dealers, by auction: it asks them to bid \ts{7}{8}{1:10}. The New York Fed's web page lists the last temporary
open market operations \ts{7}{8}{2:10}. Before the crisis the Fed was in the repo market almost every day; with excess reserves, rarely \ts{7}{8}{2:31}:
\begin{itemize}
  \item the latest is a \emph{reverse repo}: the Fed borrows money from dealers against Treasuries, agencies or MBS; dealers submitted \$4.3 billion, the Fed took
        \$2.25 billion, from Thursday the 13th to Monday the 17th, ``part of an operational readiness program'': the Fed is practising a way to drain its excess
        reserves \ts{7}{8}{2:51};
  \item for an actual \emph{repo}, which adds reserves, one must go back to August \ts{7}{8}{5:37}.
\end{itemize}
For a long time the Fed refused to call its borrowing ``reverse repos'' and called them \emph{matched sale--purchases} (MSPs): ``I'm the Fed, I print this stuff,
why do I need to provide collateral?'' Now it calls them reverses, acting like a bank \ts{7}{8}{4:56}.

\paragraph{One more level.} A Fed repo lends reserves to a dealer against a Treasury bill; the dealer has no account at the Fed, so the reserves end up at the
dealer's bank, credited to its deposit (or repaying a bank loan); if they are excess, the bank lends them in the Fed funds market. This is how open market
operations influence the Fed funds rate \ts{7}{8}{6:42}:
\begin{taccts}
\tacct[2.0cm]{Fed}{$+$Repo loan to dealer}{$+$Reserves of bank}\quad
\tacct[2.0cm]{Dealer}{$+$Deposit at bank}{$+$Repo from Fed}\quad
\tacct[2.0cm]{Dealer's bank}{$+$Reserves}{$+$Deposit of dealer}
\end{taccts}

% ------------------------------------------------------------------
\section{Interest rate spreads}\label{sec:rp-spreads}

\subsection{Before the crisis: repo $<$ Fed funds}
Stigum puzzles over why Fed funds traded above repo, and gives two reasons \ts{7}{9}{0:30}:
\begin{enumerate}
  \item \textbf{a credit spread}: repo is secured, Fed funds are not. Mehrling was never convinced, and the present reversal disproves it: nobody lends unsecured
        with any thought of default; if you fear default you cut the credit line, you do not ask a higher rate. ``Who risks a million dollars to make 137 of
        interest overnight?'' And the borrowers are banks with access to the Fed \ts{7}{9}{1:10};
  \item \textbf{segmentation}: some must borrow Fed funds and have no access to repo. But many can arbitrage, certainly any bank, by borrowing in repo and lending
        Fed funds \ts{7}{9}{2:32}.
\end{enumerate}

\paragraph{Mehrling's explanation: Bagehot's two rates.} In Bagehot's world (lecture~9) there was a \emph{market rate} of interest, made by bill brokers in the
London money market, and the official \emph{bank rate}, the Bank of England's discount rate, normally above it; nobody came to the Bank until stress pushed the market
rate up to bank rate, and then everybody flooded to it. That was monetary policy then \ts{7}{9}{2:53}. Before the crisis the US
worked the same way: the analogue of bank rate was the Fed funds target, the analogue of the market rate was repo \ts{7}{9}{4:19}. Repo below Fed funds is
\emph{discipline}: repo is the cheapest funding, so dealers repo out all their securities first; if they need more, they must go to more expensive money, borrowing
from their bank not at the Fed funds rate but at Fed funds plus 50 basis points \ts{7}{9}{4:40}.

\subsection{After the crisis: repo $>$ Fed funds}
Since January repo has moved in a band of about 16 to 30 basis points, always above Fed funds effective, sometimes by a lot \ts{7}{10}{0:09}. With the
same Bagehot logic: a market rate above the official rate means the Fed is trying to be very elastic, to pull the market rate down by pegging Fed funds below it. That
offers an arbitrage: borrow Fed funds, lend in repo, pick up about 15 basis points, and only banks (mostly) can do it \ts{7}{10}{0:51}. That
repo \emph{stays} above Fed funds despite the incentive is worrying: there are people who need liquidity and keep bidding for it, and the arbitrage is not enough to meet
the needs of the dollar funding markets, which are under stress \ts{7}{10}{1:54}. This is what the Fed would have been watching when deciding
on another QE: pegging Fed funds low was not enough, so now it buys MBS outright, \$40 billion a month, injecting reserves directly and not temporarily
\ts{7}{10}{2:36}. Will it work? Keynes said monetary policy is better at slowing the economy than at speeding it up: ``pushing on a string''
(\cref{sec:nh-managing}) \ts{7}{10}{3:37}.

\begin{finance}[Reading spreads as clues]
The method: look at patterns of interest rates and of quantities and balance sheets, and ask what must be happening behind the scenes; spreads are ``detective clues'' to
where the strains are \ts{7}{10}{4:37}. In the fall of 2007 the spread of LIBOR over Fed funds reached 100 basis points, never seen in 15 years of teaching;
knowing the two markets, it pointed to Europe \ts{7}{10}{4:58}. Prices are visible every day, balance sheets only later; ``the spreads are the key'', while the
intermediate macro textbook has just one interest rate \ts{7}{10}{5:39}. A 15 basis point reversal looks tiny, but when spreads reverse the
incentives and the flows reverse: big information in small numbers \ts{7}{10}{6:42}. Mehrling's hypothesis was formed in real time, from the FT alone, to be
tested against what happens next \ts{7}{10}{7:23}.
\end{finance}

\begin{definition}[New concepts: LIBOR]
\emph{LIBOR} (London Interbank Offered Rate) was the average rate at which a panel of banks said they could borrow unsecured from other banks in London, for each
currency and maturity; US dollar LIBOR is the Eurodollar rate (lecture~8). (After the manipulation scandal that broke in 2012, LIBOR was phased out; US dollar LIBOR
ceased in 2023, replaced by SOFR, a rate computed from overnight Treasury repo transactions.)
\end{definition}

% ------------------------------------------------------------------
\flushside
\section*{Key formulas and ideas}
\begin{keyformulas}
\begin{tabularx}{\linewidth}{@{}l L@{}}
Repo & collateralized loan organised as sale and repurchase; fails continue the loan for free\\
Balance sheets & follow the money; dealer: reverse repo (asset), repo (liability); the Fed's language is the reverse\\
Margin & always to the lender of money (hierarchy): haircut $h$, amount lent $L=(1-h)P$\\
Repo interest & $L\,r\,d/360$, simple interest, 360-day year\\
Dealers & borrow overnight (repo), lend term (reverse): banks; overnight repo $\approx$ \$1.8 tn, larger than M1\\
Liquidity spiral & haircuts up $\Rightarrow$ funding gap $\Rightarrow$ fire sales $\Rightarrow$ prices down, haircuts up\\
Fed & open market operations = repos/reverse repos with primary dealers\\
Spreads & normal: repo $<$ FF $<$ LIBOR (repo as Bagehot's market rate below bank rate); repo $>$ FF signals funding stress\\
\end{tabularx}
\end{keyformulas}

\section*{Exercises}

\begin{exercise}[Repo arithmetic]
A Treasury note with market value 102.50 per 100 face is repoed for 3 days (a weekend), haircut 2\%, repo rate 0.30\%.
\begin{enumerate}[(a)]
  \item How much is lent on \$10 million face value, and how much is repaid?
  \item What is the same interest with compound interest on a 365-day year? Compare.
\end{enumerate}
\end{exercise}

\begin{exercise}[Funding gap]
A dealer has capital 5 and holds MBS worth 100, financed by repo with a 5\% haircut.
\begin{enumerate}[(a)]
  \item Write its balance sheet.
  \item The haircut rises to 20\%. How much new funding does it need? If it must sell MBS at 90\% of value, what is its capital after selling enough to restore the
        haircut constraint?
\end{enumerate}
\end{exercise}

\begin{exercise}[Who is doing repo?]
For each statement, write the T-account entry of every party: (a) ``the Fed did \$5 billion of reverse repos with primary dealers''; (b) ``the money fund reversed in
\$1 billion of Treasuries from a dealer''; (c) ``the dealer repoed out its inventory of corporate bonds''.
\end{exercise}

\begin{exercise}[The arbitrage]
Repo trades at 0.31\%, Fed funds at 0.16\%. Describe the arbitrage a bank could do, write its T-accounts, and give two reasons why it might not close the spread.
\end{exercise}
